\documentclass[11pt]{article}
\usepackage[margin=1in]{geometry}
\usepackage{amsmath,amssymb,amsthm}
\usepackage{booktabs}
\usepackage{hyperref}

\newtheorem{proposition}{Proposition}
\newtheorem{lemma}{Lemma}
\theoremstyle{definition}
\newtheorem{definition}{Definition}

\newcommand{\R}{\mathbb{R}}
\newcommand{\Sk}{\mathcal{S}_k}
\newcommand{\dist}{\operatorname{dist}}
\newcommand{\spn}{\operatorname{span}}

\title{The quantity $\dist\{g_k/\|g_k\|,\Sk\}$ in Table~1.1 of\\
Hager \& Zhang (2013)}
\author{Companion note \texttt{1800-02} to \texttt{1800-HagerZhang\_2013}}
\date{}

\begin{document}
\maketitle

\section{What the paper says}

Table~1.1~\cite[p.~2151]{HZ13} (PALMER1C, $n=8$) reports
$\dist\{g_k,\Sk\}/\|g_k\|$, described in the text~\cite[pp.~2151--2152]{HZ13}
as the distance between $g_k/\|g_k\|$ and the space $\Sk$ spanned by the
previous search directions, up to 6 of them. The paper defines distance in
its \emph{Notation} paragraph, the last paragraph of \S1~\cite[p.~2153]{HZ13}:
\[
\dist\{x,S\}=\inf\{\|y-x\|:\ y\in S\}.
\]

\section{Definition}

\begin{definition}
Let $m_k=\min\{k,6\}$ and
$\Sk=\spn\{d_{k-1},\dots,d_{k-m_k}\}$, with $\mathcal S_0=\{0\}$. For $g_k\ne0$,
\[
\rho_k:=\dist\Bigl\{\frac{g_k}{\|g_k\|},\ \Sk\Bigr\}.
\]
\end{definition}

\begin{lemma}[Orthogonal projector]\label{lem}
Let $S\ne\{0\}$ be a subspace of $\R^n$, $Z$ a matrix with orthonormal columns
($Z^{\top}Z=I$) spanning $S$, and $P=ZZ^{\top}$. For every $v\in\R^n$:
\begin{enumerate}
\item[(i)] $Pv\in S$, and $Pv=v$ if $v\in S$;
\item[(ii)] $(I-P)v\perp S$;
\item[(iii)] $\|v\|^2=\|Pv\|^2+\|(I-P)v\|^2$ and $\|Pv\|=\|Z^{\top}v\|$;
\item[(iv)] $Pv=0\iff Z^{\top}v=0\iff v\perp S$.
\end{enumerate}
\end{lemma}

\begin{proof}
(i) $Pv=Z(Z^{\top}v)$ is a linear combination of the columns of $Z$, so
$Pv\in S$. If $v\in S$, write $v=Zu$; then $Pv=Z(Z^{\top}Z)u=Zu=v$.

(ii) Every $s\in S$ is $s=Zu$. Then
$s^{\top}(I-P)v=u^{\top}Z^{\top}v-u^{\top}(Z^{\top}Z)Z^{\top}v
=u^{\top}Z^{\top}v-u^{\top}Z^{\top}v=0$.

(iii) $v=Pv+(I-P)v$ with $Pv\in S$ by (i) and $(I-P)v\perp S$ by (ii), so the
two terms are orthogonal and Pythagoras gives the first identity. For the
second, $\|Pv\|^2=v^{\top}Z(Z^{\top}Z)Z^{\top}v=v^{\top}ZZ^{\top}v=\|Z^{\top}v\|^2$.

(iv) If $Pv=0$, then $\|Z^{\top}v\|=\|Pv\|=0$ by (iii). If $Z^{\top}v=0$, then
$Pv=Z(Z^{\top}v)=0$. Finally, $Z^{\top}v=0$ says that $v$ is orthogonal to
every column of $Z$, hence to every linear combination of them, i.e.\
$v\perp S$; and conversely $v\perp S$ implies $v$ is orthogonal to each
column, so $Z^{\top}v=0$.
\end{proof}

\begin{proposition}\label{prop}
Let $Z$ have orthonormal columns spanning $\Sk$, and $P=ZZ^{\top}$
($P=0$ if $\Sk=\{0\}$). Then
\begin{enumerate}
\item[(a)] $\displaystyle\rho_k=\frac{\dist\{g_k,\Sk\}}{\|g_k\|}$
  \quad (so the table heading and the text agree);
\item[(b)] $\displaystyle\rho_k=\frac{\|(I-P)g_k\|}{\|g_k\|}
  =\sqrt{1-\frac{\|Z^{\top}g_k\|^2}{\|g_k\|^2}}$;
\item[(c)] $0\le\rho_k\le1$, with $\rho_k=1\iff g_k\perp\Sk$ and
  $\rho_k=0\iff g_k\in\Sk$.
\end{enumerate}
\end{proposition}

\begin{proof}
Write $g=g_k\ne0$ and $S=\Sk$.

\emph{Case $S=\{0\}$.} Then $P=0$, $\dist\{g,\{0\}\}=\|g\|$ and
$\rho_k=\bigl\|g/\|g\|\bigr\|=1$, so (a)--(c) hold directly ($g\perp\{0\}$
and $g\notin\{0\}$). From now on $S\ne\{0\}$, so Lemma~\ref{lem} applies.

\emph{(a)} Since $\|g\|>0$ and $S$ is a subspace, $s\mapsto\|g\|s$ is a
bijection of $S$ onto $S$: every $y\in S$ is $y=\|g\|s$ for exactly one
$s\in S$. Hence
\[
\dist\{g,S\}=\inf_{y\in S}\|g-y\|
=\inf_{s\in S}\bigl\|g-\|g\|s\bigr\|
=\|g\|\inf_{s\in S}\Bigl\|\frac{g}{\|g\|}-s\Bigr\|
=\|g\|\,\rho_k .
\]

\emph{(b)} Fix $y\in S$ and write $g-y=(I-P)g+(Pg-y)$. The first term is
orthogonal to $S$ (Lemma~\ref{lem}(ii)); the second lies in $S$ (Lemma~\ref{lem}(i)
and $y\in S$). By Pythagoras,
\[
\|g-y\|^2=\|(I-P)g\|^2+\|Pg-y\|^2\ \ge\ \|(I-P)g\|^2 ,
\]
with equality for $y=Pg\in S$. So $\dist\{g,S\}=\|(I-P)g\|$, and (a) gives
$\rho_k=\|(I-P)g\|/\|g\|$. Dividing Lemma~\ref{lem}(iii) by $\|g\|^2$ gives
$1=\|Z^{\top}g\|^2/\|g\|^2+\rho_k^2$; since $\rho_k\ge0$, this is the second
formula.

\emph{(c)} By (b) and Lemma~\ref{lem}(iii),
$\rho_k^2=1-\|Pg\|^2/\|g\|^2$ with $0\le\|Pg\|^2\le\|g\|^2$, so
$0\le\rho_k\le1$. Moreover
\begin{itemize}
\item $\rho_k=1\iff Pg=0\iff g\perp S$, by Lemma~\ref{lem}(iv);
\item $\rho_k=0\iff (I-P)g=0\iff g=Pg$. If $g=Pg$ then $g\in S$ by
  Lemma~\ref{lem}(i); if $g\in S$ then $Pg=g$, also by Lemma~\ref{lem}(i).
\qedhere
\end{itemize}
\end{proof}

In exact CG on a quadratic, $g_k\perp d_j$ for all $j<k$
\cite[pp.~2159--2160]{HZ13}, so $\rho_k=1$ always. Any value below $1$ in
Table~1.1 is therefore lost orthogonality. The paper's test (3.2) is simply
$\rho_k\le\eta_0$.

\section{Two examples}

Take $n=3$, $A=\operatorname{diag}(1,2,3)$, $g_0=(1,1,1)^{\top}$, and the
Fletcher--Reeves code of \S1 of the paper~\cite[p.~2152]{HZ13}:
\[
d_0=-g_0,\quad \alpha_k=\frac{\|g_k\|^2}{d_k^{\top}Ad_k},\quad
g_{k+1}=g_k+\alpha_kAd_k,\quad
\beta_k=\frac{\|g_{k+1}\|^2}{\|g_k\|^2},\quad
d_{k+1}=-g_{k+1}+\beta_kd_k .
\]
Here $k\le3<6$, so $\Sk=\spn\{d_0,\dots,d_{k-1}\}$.

\subsection*{Example 1: exact CG (orthogonality kept)}

\emph{Step $k=0$.} $Ad_0=(-1,-2,-3)^{\top}$, $d_0^{\top}Ad_0=1+2+3=6$,
$\|g_0\|^2=3$, so $\alpha_0=\tfrac36=\tfrac12$ and
\begin{align*}
g_1&=(1,1,1)+\tfrac12(-1,-2,-3)=(\tfrac12,\,0,\,-\tfrac12),
  \qquad \|g_1\|^2=\tfrac12,\qquad \beta_0=\tfrac{1/2}{3}=\tfrac16,\\
d_1&=-(\tfrac12,0,-\tfrac12)+\tfrac16(-1,-1,-1)=(-\tfrac23,\,-\tfrac16,\,\tfrac13).
\end{align*}

\emph{Step $k=1$.} $Ad_1=(-\tfrac23,-\tfrac13,1)^{\top}$,
$d_1^{\top}Ad_1=\tfrac49+\tfrac1{18}+\tfrac13=\tfrac56$, so
$\alpha_1=\tfrac{1/2}{5/6}=\tfrac35$ and
\begin{align*}
g_2&=(\tfrac12,0,-\tfrac12)+\tfrac35(-\tfrac23,-\tfrac13,1)
   =(\tfrac1{10},\,-\tfrac15,\,\tfrac1{10}),
  \qquad \|g_2\|^2=\tfrac3{50},\qquad \beta_1=\tfrac{3/50}{1/2}=\tfrac3{25},\\
d_2&=-(\tfrac1{10},-\tfrac15,\tfrac1{10})+\tfrac3{25}(-\tfrac23,-\tfrac16,\tfrac13)
   =(-\tfrac9{50},\,\tfrac9{50},\,-\tfrac3{50}).
\end{align*}

\emph{Step $k=2$.} $Ad_2=(-\tfrac9{50},\tfrac9{25},-\tfrac9{50})^{\top}$,
$d_2^{\top}Ad_2=\tfrac{81+162+27}{2500}=\tfrac{27}{250}$, so
$\alpha_2=\tfrac{3/50}{27/250}=\tfrac59$ and
\[
g_3=(\tfrac1{10},-\tfrac15,\tfrac1{10})+\tfrac59(-\tfrac9{50},\tfrac9{25},-\tfrac9{50})
=(0,\,0,\,0).
\]

\emph{Values of $\rho_k$.} $g_1^{\top}d_0=-(\tfrac12+0-\tfrac12)=0$, so
$g_1\perp\mathcal S_1$ and $\rho_1=1$ (Proposition~\ref{prop}(c)).
$g_2^{\top}d_0=-(\tfrac1{10}-\tfrac15+\tfrac1{10})=0$ and
$g_2^{\top}d_1=-\tfrac1{15}+\tfrac1{30}+\tfrac1{30}=0$, so $\rho_2=1$.
At $k=3$, $g_3=0$: the minimum is reached in $n$ steps.

\subsection*{Example 2: one wrong step (orthogonality lost)}

Same problem, but the first step length is $\tilde\alpha_0=\tfrac35$ instead
of $\tfrac12$ (standing in for a rounding error). All later steps follow the
code exactly.

\emph{Step $k=0$.} With $Ad_0=(-1,-2,-3)^{\top}$ as before,
\begin{align*}
g_1&=(1,1,1)+\tfrac35(-1,-2,-3)=(\tfrac25,\,-\tfrac15,\,-\tfrac45),
  \qquad \|g_1\|^2=\tfrac{21}{25},\qquad \beta_0=\tfrac{21/25}{3}=\tfrac7{25},\\
d_1&=-(\tfrac25,-\tfrac15,-\tfrac45)+\tfrac7{25}(-1,-1,-1)
   =(-\tfrac{17}{25},\,-\tfrac2{25},\,\tfrac{13}{25}).
\end{align*}

\emph{Step $k=1$.} $Ad_1=(-\tfrac{17}{25},-\tfrac4{25},\tfrac{39}{25})^{\top}$,
$d_1^{\top}Ad_1=\tfrac{289+8+507}{625}=\tfrac{804}{625}$, so
$\alpha_1=\tfrac{21/25}{804/625}=\tfrac{175}{268}$ and
\[
g_2=\bigl(\tfrac25-\tfrac{119}{268},\ -\tfrac15-\tfrac{28}{268},\
-\tfrac45+\tfrac{273}{268}\bigr)
=\tfrac1{1340}(-59,\,-408,\,293).
\]

\emph{Step $k=2$.} The fractions become large; numerically,
$g_3\approx(-0.0891,\ 0.0632,\ -0.0287)\ne0$: the method no longer finishes
in $n$ steps.

\emph{Value of $\rho_1$.} $\mathcal S_1=\spn\{(1,1,1)^{\top}\}$, with
$Z=\tfrac1{\sqrt3}(1,1,1)^{\top}$. Then $Pg_1=-\tfrac15(1,1,1)^{\top}$ and
$(I-P)g_1=(\tfrac35,0,-\tfrac35)^{\top}$, so by Proposition~\ref{prop}(b)
\[
\rho_1=\frac{3\sqrt2/5}{\sqrt{21}/5}=\sqrt{\tfrac67}\approx0.926 .
\]

\emph{Value of $\rho_2$.} $\rho_2$ depends only on the direction of $g_2$ and
on $\mathcal S_2$, so we may use $g=(-59,-408,293)^{\top}$ and the spanning
vectors $(1,1,1)^{\top}\propto d_0$, $(17,2,-13)^{\top}\propto d_1$.
Gram--Schmidt gives the orthonormal basis
\[
z_1=\tfrac1{\sqrt3}(1,1,1)^{\top},\qquad
(17,2,-13)-\tfrac{17+2-13}{3}(1,1,1)=(15,0,-15)
\ \Rightarrow\ z_2=\tfrac1{\sqrt2}(1,0,-1)^{\top}.
\]
Then $z_1^{\top}g=-\tfrac{174}{\sqrt3}$, $z_2^{\top}g=-\tfrac{352}{\sqrt2}$,
$\|Z^{\top}g\|^2=10092+61952=72044$ and $\|g\|^2=255794$. By
Proposition~\ref{prop}(b),
\[
\rho_2=\sqrt{1-\tfrac{72044}{255794}}=\sqrt{\tfrac{13125}{18271}}\approx0.848 .
\]

\begin{center}
\begin{tabular}{@{}ccc@{}}
\toprule
$k$ & $\rho_k$ (exact) & $\rho_k$ (wrong step)\\
\midrule
0 & 1 & 1\\
1 & 1 & 0.926\\
2 & 1 & 0.848\\
\bottomrule
\end{tabular}
\end{center}
\section{Reading Table~1.1}

\begin{itemize}
\item L-BFGS never drops below $0.64$: orthogonality is essentially kept.
\item CG at $k=5$: $\rho_5\approx3\times10^{-9}$, so $g_5$ essentially lies
  in $\spn\{d_0,\dots,d_4\}$.
\item For $k>6$ only the last 6 directions are kept, which is why CG can jump
  back near $1$ (e.g.\ $k=9$).
\end{itemize}

\begin{thebibliography}{9}
\bibitem{HZ13} W.~W.~Hager and H.~Zhang, The limited memory conjugate gradient
method, \emph{SIAM J. Optim.}, 23 (2013), pp.~2150--2168.
\end{thebibliography}

\end{document}
